\section*{Bab \ref{ch:b3:holomorphic} --- Fungsi
Holomorfik}

\begin{solution}{exo:b3:holomorphic:1}
(a) Untuk $\bar z$: $P = x$, $Q = -y$, sehingga $P_x = 1 \neq -1 = Q_y$:
jadi tak dapat diturunkan-$\C$ di mana pun. Untuk $\abs z^2$: $P = x^2 + y^2$, $Q =
0$: sehingga Cauchy--Riemann menuntut $2x = 0 = 2y$: jadi dapat diturunkan di
$0$ saja --- dan \omterm{def:b3:holomorphic:holo}{holomorfik} di mana pun tidak (sebab tak ada \omterm{def:b3:topology:topology}{himpunan terbuka}).
Untuk $\operatorname{Re}z$: $P_x = 1 \neq 0 = Q_y$: jadi tak di mana pun.

(b) Berlaku $x^2 - y^2 = \operatorname{Re}(z^2)$: sehingga $f(z) = z^2$ berhasil.
Untuk semua solusinya: jika $\operatorname{Re}f = \operatorname{Re}g$
dengan $f, g$ \omterm{def:b3:holomorphic:holo}{holomorfik} pada $\C$ yang \omterm{def:b3:topology:connected}{terhubung}, maka $h = f -
g$ mempunyai $\operatorname{Re}h = 0$; lalu Cauchy--Riemann memberikan $h' =
P_x + \iu Q_x = 0 - \iu P_y = 0$: sehingga $h$ merupakan konstanta
imajiner. Jawabannya: $f(z) = z^2 + \iu c$ dengan $c \in \R$.
\end{solution}

\begin{solution}{exo:b3:holomorphic:2}
Pada lingkaran satuannya $\gamma(t) = \eu^{\iu t}$:
\[
\int_C z^n\,\dd z = \int_0^{2\pi}\eu^{\iu nt}\,\iu\eu^{\iu
t}\dd t = \iu\int_0^{2\pi}\eu^{\iu(n+1)t}\dd t
= \begin{cases} 2\iu\pi & n = -1,\\ 0 & n \neq -1.\end{cases}
\]
Untuk $\bar z$: sepanjang $[0, 1+\iu]$ dengan $\gamma(t) = t(1 + \iu)$:
$\int_0^1 t(1 - \iu)(1 + \iu)\dd t = \int_0^12t\,\dd t = 1$.
Sedangkan sepanjang $0 \to 1 \to 1 + \iu$: $\int_0^1t\dd t +
\int_0^1(1 - \iu t)\,\iu\,\dd t = \frac12 + \iu + \frac12 = 1
+ \iu$. Jadi nilainya berbeda antara ujung yang sama: sehingga menurut
\cref{prop:b3:holomorphic:primitive}, $\bar z$ tak berantiturunan
pada \omterm{def:b3:topology:topology}{himpunan terbuka} mana pun yang memuat kedua \omterm{def:b3:topology:pathconnected}{lintasannya}.
\end{solution}

\begin{solution}{exo:b3:holomorphic:3}
(a) Bidang teririsnya $\Omega = \C\setminus\intoc{-\infty}0$ bersifat
berbentuk bintang terhadap $1$, dan $\frac1z \in \mathcal H(\Omega)$:
sehingga \cref{thm:b3:holomorphic:cauchy} menyediakan antiturunan $L$
dengan $L(1) = 0$. Lalu $\bigl(z\eu^{-L(z)}\bigr)' =
\eu^{-L}(1 - z\cdot\frac1z) = 0$: sehingga $z = c\,\eu^{L(z)}$ dengan $c
= 1$ (di $z = 1$). Dengan menulis $L = u + \iu v$: $\abs z = \eu^u$
dan $z = \abs z\eu^{\iu v}$ dengan $v$ yang \omterm{def:b3:topology:continuity}{kontinu}, $v(1) = 0$, dan
$v \in \intoo{-\pi}\pi$ (sebab $v$ merupakan argumen \omterm{def:b3:topology:continuity}{kontinu} $z$
pada $\Omega$ yang \omterm{def:b3:topology:connected}{terhubung}, sehingga petanya menghindari kelipatan
gasal $\pi$ --- karena tak ada titik $\Omega$ yang terletak pada $\R_-$
--- dan, karena memuat $v(1) = 0$, ia tinggal di $\intoo{-\pi}\pi$:
jadi $v$ adalah argumen pokoknya): sehingga $L = \log$.

(b) Seandainya $g$ merupakan logaritma \omterm{def:b3:topology:continuity}{kontinu} pada $\C^*$: maka $h(t) =
g(\eu^{\iu t})$ memenuhi $\eu^{h(t)} = \eu^{\iu t}$, sehingga $h(t)
- \iu t \in 2\iu\pi\Z$, dan lewat \omterm{def:b3:topology:continuity}{kekontinuannya} $h(t) = \iu t +
2\iu\pi k$ untuk sebuah bilangan bulat $k$ yang tetap. Lalu $g(1) = h(0) =
2\iu\pi k$ dan $g(1) = h(2\pi) = 2\iu\pi(k + 1)$: yang
bertentangan.

(c) Pada $D(0,1)$: $\log(1 + z) =
\sum_{n\geq1}\frac{(-1)^{n+1}}{n}z^n$ --- sebab kedua ruasnya lenyap di
$0$ dan berturunan $\frac1{1+z} = \sum(-1)^nz^n$
(\cref{ex:b3:holomorphic:examples}); dan antiturunan bersifat tunggal sampai
sebuah konstanta pada cakram yang \omterm{def:b3:topology:connected}{terhubung}.
\end{solution}

\begin{solution}{exo:b3:holomorphic:4}
(a) Uraikan di $0$ (dengan jari-jari $\infty$): lalu menurut \omterm{thm:b3:holomorphic:analytic}{taksiran Cauchy}
pada $C_r$, $\abs{c_k} \leq C(1 + r)^n/r^k \to 0$ saat $r \to
\infty$ untuk $k > n$: sehingga $f = \sum_{k\leq n}c_kz^k$.

(b) Jika $\operatorname{Re}f \leq M$: maka $g = \eu^f$ utuh dengan
$\abs g = \eu^{\operatorname{Re}f} \leq \eu^M$: jadi konstan menurut
Liouville. Lalu $g' = f'g = 0$ dengan $g$ yang tak lenyap: sehingga $f' =
0$, dan $f$ konstan.

(c) Jika $f$ melewatkan setengah bidang $H$, maka sebuah pemetaan afin $w \mapsto
\alpha w + \beta$ mengirim $\C\setminus H$ ke dalam
$\{\operatorname{Re} \leq M\}$; lalu terapkan (b) pada $\alpha f +
\beta$.
\end{solution}

\begin{solution}{exo:b3:holomorphic:5}
(a) Fungsi $g(z) = f(z) - z^2$ lenyap di titik $\frac1n$,
yang menumpuk di $0 \in \Omega$: sehingga menurut teorema keidentikannya
(karena $\Omega$ \omterm{def:b3:topology:connected}{terhubung}), $g \equiv 0$: jadi $f(z) = z^2$.

(b) Ya: $f(z) = z\cos(\pi/z)$ bersifat \omterm{def:b3:holomorphic:holo}{holomorfik} pada $\C^*$
(lewat komposisinya) dan $f(\frac1n) = \frac1n\cos(n\pi) =
\frac{(-1)^n}n$. Tak ada pertentangan dengan (a): sebab titik tumpuk
$0$ milik simpul interpolasinya \emph{tak} termasuk
$\C^*$, sehingga teorema keidentikannya bungkam --- jadi dua fungsi
yang berbeda ($z\cos(\pi/z)$ dan, katakanlah, hasil interpolasi
lain) boleh berbagi nilai ini.

(c) Ambil $f \equiv 0$ di mana-mana, terhadap $g = 0$ pada $D(0,1)$ dan
$g = 1$ pada $D(3,1)$: keduanya \omterm{def:b3:holomorphic:holo}{holomorfik} pada gabungan yang tak \omterm{def:b3:topology:connected}{terhubung},
sama pada $D(0,1)$, tetapi berbeda. Sebab argumen terbuka-tertutup
teorema keidentikannya memerlukan keterhubungan untuk merambat dari satu
komponen ke komponen lainnya --- dan ia tak dapat.
\end{solution}

\begin{solution}{exo:b3:holomorphic:6}
(a) Menurut asas maksimum yang diterapkan pada daerah terbatasnya:
$\sup_{\mathbb D}\abs f = \sup_{\partial\mathbb D}\abs f = 1$.
Karena $f$ tak berakar, $1/f$ bersifat \omterm{def:b3:holomorphic:holo}{holomorfik} pada $\mathbb D$,
\omterm{def:b3:topology:continuity}{kontinu} pada \omterm{def:b3:topology:interior}{penutupnya}, dengan modulus batas $1$: sehingga demikian pula
$\abs{1/f} \leq 1$, yakni $\abs f \geq 1$. Jadi $\abs f \equiv
1$: sehingga modulusnya mencapai maksimum di \omterm{def:b3:topology:interior}{interiornya}, dan
\cref{thm:b3:holomorphic:maximum}(2) memaksa $f$ konstan.

(b) Ambil $f(z) = z$: dengan modulus batas $1$, akar di titik asalnya, dan
tak konstan --- sebab akarnya tepat merupakan hal yang menghalangi argumen
$1/f$-nya.
\end{solution}

\begin{solution}{exo:b3:holomorphic:7}
\omterm{def:b3:topology:continuity}{Kekontinuan} $G$: lewat kekonvergenan terdominasi dengan pendominasi $h$
(yakni batas seragam lokalnya). \omterm{def:b3:holomorphic:holo}{Keholomorfikannya} lewat Morera (yang ditegakkan
di dalam \cref{thm:b3:holomorphic:weierstrassconv}): untuk sebuah
segitiga tertutup $T$ di sebuah cakram yang $\abs g \leq h$-nya,
\[
\int_{\partial T}G(z)\,\dd z =
\int_X\Bigl(\int_{\partial T}g(x, z)\,\dd z\Bigr)\dd\mu(x) =
0,
\]
dengan penukarannya lewat Fubini (sebab $\int_X\int_{\partial T}\abs
g \leq \operatorname{length}(\partial T)\int h < \infty$) dan
kelenyapan bagian dalamnya lewat Goursat. Untuk $\Gamma$: pada jalur $a
\leq \operatorname{Re}z \leq b$ (dengan $0 < a \leq b$),
$\abs{t^{z-1}\eu^{-t}} = t^{\operatorname{Re}z-1}\eu^{-t} \leq
(t^{a-1} + t^{b-1})\eu^{-t}$, yang \omterm{def:b3:lebesgue:l1}{terintegralkan} pada
$\intoo0{+\infty}$: sehingga $\Gamma$ \omterm{def:b3:holomorphic:holo}{holomorfik} pada
$\{\operatorname{Re} z > 0\}$ (sebab \omterm{def:b3:measure:measure}{ukurannya} hanya
berhingga-$\sigma$, tetapi argumennya hanya memerlukan pendominasi
yang \omterm{def:b3:lebesgue:l1}{terintegralkan}). Lalu menurut teorema keidentikannya, persamaan fungsional
$\Gamma(z + 1) = z\Gamma(z)$, yang dibuktikan pada
$\intoo0{+\infty}$ (\cref{ex:b3:lebesgue:gamma}), berlaku pada
seluruh setengah bidangnya.
\end{solution}

\begin{solution}{exo:b3:holomorphic:8}
Tulis $P = c\prod_k(X - a_k)^{m_k}$
(\cref{pb:b3:holomorphic:1}). Misalkan $P'(z) = 0$. Jika $P(z) = 0$,
maka $z$ salah satu dari $a_k$: jadi di selubungnya. Selain itu,
turunan logaritmiknya memberikan
\[
0 = \frac{P'(z)}{P(z)} = \sum_k\frac{m_k}{z - a_k}
= \sum_k m_k\,\frac{\bar z - \bar a_k}{\abs{z - a_k}^2} ;
\]
lalu dengan mengonjugatkannya, $\sum_kw_k(z - a_k) = 0$ dengan $w_k = m_k/\abs{z -
a_k}^2 > 0$: sehingga $z = \sum_k\frac{w_k}{\sum w}\,a_k$, yakni kombinasi
cembung akarnya. Untuk $P = z^3 - 1$: akarnya adalah akar pangkat tiga
satuan, dan $P' = 3z^2$ berakar rangkap $0$ --- yakni
sentroid segitiga sama sisinya.
\end{solution}

\begin{solution}{exo:b3:holomorphic:9}
Bujur sangkar satuan tertutup $K = \{x + \iu y : 0 \leq x, y \leq
1\}$ bersifat \omterm{def:b3:topology:compact}{kompak}: sehingga $M = \sup_K\abs f < \infty$. Setiap $z \in \C$
berselisih dari sebuah titik $K$ oleh sebuah unsur $\Z + \iu\Z$
(kurangkan bagian bulatnya), dan $f$ invarian terhadap
translasi itu (iterasikan kedua hubungannya): sehingga $\abs f \leq M$ pada
$\C$. Lalu Liouville: $f$ konstan. Karena itu sembarang fungsi meromorfik
berperiode ganda yang tak konstan --- yakni fungsi eliptik
teori klasiknya --- haruslah berkutub.
\end{solution}

\begin{solution}{exo:b3:holomorphic:10}
(a) Ambillah bagian real pada rumus nilai rata-ratanya
(\cref{thm:b3:holomorphic:maximum}(1)). Untuk kemulusannya: $f$ bersifat
analitik, sehingga $P, Q \in \mathcal C^\infty$; lalu dengan menurunkan
Cauchy--Riemann: $P_{xx} = (Q_y)_x = (Q_x)_y = (-P_y)_y =
-P_{yy}$ (lewat kesetangkupan Schwarz turunan keduanya): sehingga $\Delta P =
0$.

(b) Jika $\operatorname{Re}f$ mencapai maksimum di \omterm{def:b3:topology:interior}{interior} sebuah
$\Omega$ yang \omterm{def:b3:topology:connected}{terhubung}: maka $g = \eu^f$ mempunyai $\abs g =
\eu^{\operatorname{Re}f}$ yang mencapai maksimum di \omterm{def:b3:topology:interior}{interiornya}, sehingga
$g$, jadi $\operatorname{Re}f = \ln\abs g$, bersifat konstan
(\cref{thm:b3:holomorphic:maximum}(2)). Sedangkan pada daerah terbatas
dengan \omterm{def:b3:topology:continuity}{kekontinuan} sampai batasnya,
$\sup_{\bar\Omega}\operatorname{Re}f =
\sup_{\partial\Omega}\operatorname{Re}f$.
\end{solution}

\begin{solution}{exo:b3:holomorphic:11}
(a) Kedua rumusnya bersesuaian pada $I$ (sebab $z = \bar z$ dan $f$ real
di sana: $\overline{f(\bar z)} = \overline{f(z)} = f(z)$), dan
$z \mapsto \overline{f(\bar z)}$ \omterm{def:b3:topology:continuity}{kontinu} pada
$\Omega^-\cup I$ sebagai komposisi \omterm{def:b3:topology:continuity}{pemetaan kontinu}: sehingga $F$
\omterm{def:b3:topology:continuity}{kontinu} pada $\Omega$. Untuk \omterm{def:b3:holomorphic:holo}{keholomorfikannya} pada $\Omega^-$: di dekat $z_0
\in \Omega^-$, uraikan $f(w) = \sum c_n(w - \bar z_0)^n$ di dekat
$\bar z_0 \in \Omega^+$; maka
\[
\overline{f(\bar z)} = \sum_n\bar c_n\,(z - z_0)^n,
\]
yakni sebuah deret pangkat yang konvergen: jadi \omterm{def:b3:holomorphic:holo}{holomorfik}.

(b) Segitiga yang menghindari $I$ ditangani Goursat di
$\Omega^\pm$. Sedangkan untuk segitiga yang memotong $I$, irislah ia oleh sumbu
realnya menjadi paling banyak tiga segitiga atau segi empat, masing-masing bersisi
satu pada $I$; lalu untuk kepingan $P$ semacam itu yang termuat, katakanlah, di
$\overline{\Omega^+}$, \omterm{def:b3:holomorphic:contour}{integral konturnya} adalah limit saat
$\varepsilon \to 0^+$ dari integral atas translasi bertipe $P + \iu
\varepsilon$ (sebab \omterm{def:b3:topology:continuity}{kekontinuan} seragam $F$ pada
kepingan \omterm{def:b3:topology:compact}{kompaknya} membuat integral batasnya konvergen, dengan
sisi pada $I$-nya didekati dari atas), dan setiap translasinya
terletak di $\Omega^+$ tempat Goursat memberikan $0$. Lalu menjumlahkan
kepingannya: $\int_{\partial T}F = 0$. Lalu Morera (yakni kriteria
di dalam \cref{thm:b3:holomorphic:weierstrassconv}): sehingga $F$
\omterm{def:b3:holomorphic:holo}{holomorfik} pada $\Omega$.

(c) Pada cakramnya, $G(z) = \overline{f(\bar z)}$ bersifat
\omterm{def:b3:holomorphic:holo}{holomorfik} menurut perhitungan (a) dan bersesuaian dengan $f$ pada
diameternya, yakni himpunan yang bertitik tumpuk: sehingga $G = f$ di mana-mana
(menurut teorema keidentikannya). Seandainya $f$ real pada sebuah \omterm{def:b3:topology:topology}{himpunan terbuka}
$U$ yang tak kosong: maka pada $U$ kedua turunan parsial $Q = \operatorname{Im}f$
lenyap, dan Cauchy--Riemann memindahkannya ke $P =
\operatorname{Re}f$ (sebab $P_x = Q_y = 0$ dan $P_y = -Q_x = 0$), sehingga
$f' = P_x + \iu Q_x = 0$ pada $U$: jadi $f$ konstan pada $U$,
sehingga di mana-mana menurut teorema keidentikannya (karena $\Omega$
\omterm{def:b3:topology:connected}{terhubung}).
\end{solution}

\begin{solution}{exo:b3:holomorphic:12}
(a) Berlaku $1 = f^2 + g^2 = (f + \iu g)(f - \iu g)$, sehingga $h = f +
\iu g$ tak pernah lenyap (sebab kofaktornya harus meledak).
Untuk $h$ yang utuh dan bebas akar: $h'/h$ bersifat utuh, dan $\C$
berbentuk bintang, sehingga ia berantiturunan $\varphi_0$
(\cref{thm:b3:holomorphic:cauchy}); lalu
$\bigl(h\eu^{-\varphi_0}\bigr)' = \eu^{-\varphi_0}(h' -
h\varphi_0') = 0$: sehingga $h = c\,\eu^{\varphi_0}$ dengan $c \neq 0$,
dan dengan menyerap sebuah konstanta $\log c$ ke dalam $\varphi = \varphi_0
+ \log c$ (yakni sembarang logaritma kompleks $c$): $h =
\eu^{\varphi}$.

(b) Dengan $h = \eu^{\varphi}$ dan $h^{-1} = f - \iu g =
\eu^{-\varphi}$:
\[
f = \frac{\eu^{\varphi} + \eu^{-\varphi}}2,
\qquad
g = \frac{\eu^{\varphi} - \eu^{-\varphi}}{2\iu} .
\]
Dengan menulis $\varphi = \iu\psi$ dengan $\psi = -\iu\varphi$
yang utuh, keduanya berbunyi $f = \cos\psi$ dan $g = \sin\psi$: sehingga
solusi utuhnya tepat berupa pasangan $(\cos\psi,
\sin\psi)$ dengan $\psi$ yang utuh, dan konversnya adalah
identitas $\cos^2 + \sin^2 = 1$.
Untuk $f^2 + g^2 = 0$: $(f + \iu g)(f - \iu g) = 0$ pada
daerah integral $\mathcal H(\C)$ (sebab $\C$ \omterm{def:b3:topology:connected}{terhubung}: karena pembagi
nol akan melanggar teorema keidentikannya): sehingga $g = \pm\iu
f$ dengan $f$ yang utuh sembarang.
\end{solution}

\begin{solution}{pb:b3:holomorphic:1}
\textbf{1.} Untuk $\abs z \geq 1$:
\[
\abs{P(z)} \geq \abs z^n\Bigl(1 -
\frac{\abs{a_{n-1}}}{\abs z} - \dots -
\frac{\abs{a_0}}{\abs z^n}\Bigr)
\geq \abs z^n\Bigl(1 - \frac{A}{\abs z}\Bigr),
\qquad A = \sum_k\abs{a_k} :
\]
untuk $\abs z \geq R_0 = \max(1, 2A)$: $\abs{P(z)} \geq
\frac12\abs z^n \to \infty$.

\textbf{2.} Pilihlah $R \geq R_0$ dengan $\frac12R^n \geq
\abs{P(0)}$. Pada $\bar D(0, R)$ yang \omterm{def:b3:topology:compact}{kompak}, $\abs P$ yang \omterm{def:b3:topology:continuity}{kontinu}
mencapai minimum, di suatu $z_0$; sedangkan di luarnya, $\abs P
\geq \frac12R^n \geq \abs{P(0)} \geq \abs{P(z_0)}$: sehingga
minimumnya global.

\textbf{3.} Fungsi $Q(h) = P(z_0 + h)/P(z_0)$ merupakan polinomial dalam $h$
dengan $Q(0) = 1$; dan ia tak konstan (sebab $P$ demikian), sehingga suatu
koefisien di luar konstantanya tak nol: jadi $Q(h) = 1 +
c_kh^k + h^{k+1}S(h)$ dengan $k \geq 1$ yang minimal, $c_k \neq 0$, dan
$S \in \C[X]$.

\textbf{4.} Untuk akarnya: sembarang $w = \rho\eu^{\iu\varphi} \neq 0$ mempunyai
akar ke-$k$ $\rho^{1/k}\eu^{\iu\varphi/k}$, dengan
$\rho^{1/k}$ yang ada menurut teorema nilai antara yang diterapkan
pada $t \mapsto t^k$ di $\intco0\infty$ --- jadi tanpa kelingkaran.
Lalu pilihlah $\omega$ dengan $\omega^k = -1/c_k$. Maka
\[
Q(t\omega) = 1 - t^k + t^{k+1}\,\omega^{k+1}S(t\omega),
\qquad
\abs{Q(t\omega)} \leq 1 - t^k + C\,t^{k+1}
\quad (0 < t \leq 1),
\]
dengan $C = \abs\omega^{k+1}\sup_{\abs h \leq
\abs\omega}\abs{S(h)}$ (perhatikan $1 - t^k \geq 0$ pada $\intcc01$).

\textbf{5.} Untuk $0 < t < \min(1, 1/C)$: $\abs{Q(t\omega)}
\leq 1 - t^k(1 - Ct) < 1$, yakni $\abs{P(z_0 + t\omega)} <
\abs{P(z_0)}$ --- yang bertentangan dengan keminimalan globalnya. Jadi
$P(z_0) = 0$: sehingga bukti d'Alembert--Argand \omterm{def:b3:complete:complete}{lengkap}.

\textbf{6.} Untuk versi Liouvillenya: jika $P$ tak pernah lenyap, maka $1/P$
bersifat utuh; lalu menurut pertanyaan 1, $\abs{1/P} \leq 2R_0^{-n}$ di luar
$\bar D(0, R_0)$, dan $1/P$ \omterm{def:b3:topology:continuity}{kontinu} pada cakram \omterm{def:b3:topology:compact}{kompak}
itu, sehingga terbatas di sana pula: jadi utuh dan terbatas, sehingga konstan
(\cref{cor:b3:holomorphic:liouville}), yang membuat $P$ konstan:
mustahil. Bahannya: d'Alembert memakai kekompakan (yakni keberadaan
minimumnya) dan keberadaan akar ke-$k$ lewat bentuk kutubnya;
sedangkan Liouville memakai seluruh perkakas Cauchy (yakni Goursat ---
yang sendirinya argumen \omterm{def:b3:topology:compact}{kompak} bersarang --- dan \omterm{thm:b3:holomorphic:analytic}{taksiran
Cauchynya}) ditambah kekoersifan yang sama. Jadi kekompakan cakram
tertutup adalah inti bersamanya yang \omterm{def:b3:rings:divisibility}{tak tereduksi}.

\textbf{7.} Jika $\deg P \geq 1$, pilihlah sebuah akar $\alpha$
(pertanyaan 5); lalu bagilah: $P = (X - \alpha)Q + P(\alpha) = (X -
\alpha)Q$ dengan $\deg Q = n - 1$; lalu induksikan. Lalu dengan mengelompokkan faktor
yang sama: $P = c\prod_i(X - \alpha_i)^{m_i}$ dengan $\sum m_i = n$.

\textbf{8.} Untuk $P$ yang real: $P(\bar\alpha) =
\overline{P(\alpha)} = 0$, dan kelipatannya bersesuaian
(konjugatkan pemfaktorannya): sehingga akar tak realnya muncul berpasangan,
yang menyumbang $(X - \alpha)(X - \bar\alpha) = X^2 -
2\operatorname{Re}(\alpha)X + \abs\alpha^2$, yakni kuadratik real
berdiskriminan $< 0$. Karena itu diperoleh daftar \omterm{def:b3:rings:divisibility}{tak
tereduksinya}, dan polinomial real berderajat gasal, yang berakar
tak real dalam jumlah yang genap, haruslah berakar real. Untuk bukti
langsungnya: $P(x) \to \pm\infty$ saat $x \to \pm\infty$ (sebab berderajat gasal
dengan koefisien utama positif katakanlah), sehingga $P$ berganti tanda, dan
teorema nilai antaranya berlaku. Untuk $X^3 - X - 1$:
kedua argumennya memberikan satu akar real $\approx 1.3247$
(beserta sepasang sekawan).

\textbf{9.} (i) Polinomial $\chi_u \in \C[X]$ tak konstan: sehingga ia
berakar $\lambda$, dan $\det(u - \lambda\,\mathrm{id}) = 0$
memberikan sebuah vektor eigen; lalu mesin \omterm{thm:b3:modules:structure}{pembagi elementer}
\cref{thm:b3:modules:jordan} berlaku bagi sembarang matriks kompleks,
sebab $\chi$ selalu terbelah. (ii) Misalkan $P \in \bar\Q[X]$
tak konstan. Sebagai polinomial atas $\C$ ia berakar $z
\in \C$; lalu $z$ bersifat aljabar atas $\bar\Q$, sehingga atas $\Q$ lewat
ketransitifannya (\cref{cor:b3:galois:algclosed}), jadi $z \in
\bar\Q$: sehingga setiap polinomial tak konstan atas $\bar\Q$ berakar
di $\bar\Q$.

\textbf{10.} Atas $\Q(\iu)$, pertanyaan 4 sudah gagal:
sebab akar ke-$k$ tak harus ada (tak ada $\sqrt2$), dan bahkan dengan
akarnya diberikan, pertanyaan 2 gagal --- sebab barisan pemberi minimumnya tak harus
konvergen, karena $\Q$ tak \omterm{def:b3:complete:complete}{lengkap}; sedangkan pada jalur Liouville,
segitiga \omterm{def:b3:topology:compact}{kompak} bersarang Goursat beririsan kosong
atas titik $\Q(\iu)$. Jadi kedua buktinya memakan \omterm{def:b3:complete:complete}{kelengkapan}
(setara dengan itu, lewat kekonvergenan monoton terbatas, \omterm{def:b3:complete:complete}{kelengkapan}
urutan) $\R$; sedangkan keterhubungannya menggerakkan teorema
nilai antara di balik bentuk kutubnya. Jadi pernyataan
``$\C$ \omterm{def:b3:galois:closure}{tertutup secara aljabar}'' adalah aljabar; sedangkan setiap bukti
yang dikenal atasnya adalah analisis yang diselundupkan lewat definisi
$\R$.

\textbf{11.} Menurut \cref{thm:b3:holomorphic:analytic} di $a$,
$f^{(n)}(a) = n!\,c_n$ dengan
\[
c_n = \frac1{2\iu\pi}\int_{C_r}\frac{f(w)}{(w -
a)^{n+1}}\,\dd w,
\qquad
\abs{c_n} \leq \frac{2\pi r}{2\pi}\cdot
\frac{\sup_{C_r}\abs f}{r^{n+1}}
= \frac{\sup_{C_r}\abs f}{r^{n}} :
\]
yakni \omterm{thm:b3:holomorphic:analytic}{taksiran Cauchy}, dalam bentuk yang ditampilkan setelah
dikalikan $n!$. Jika $\abs f \leq M$ pada $\C$: maka untuk setiap
$a$ dan setiap $r$, $\abs{f'(a)} \leq M/r \to 0$ saat $r \to
\infty$, sehingga $f' \equiv 0$ dan $f$ konstan pada
$\C$ yang \omterm{def:b3:topology:connected}{terhubung} --- yakni Liouville yang dipulihkan.

\textbf{12.} Uraikan $f = \sum_kc_kz^k$ di $0$ (dengan jari-jari
$\infty$). Untuk $k > m$: $\abs{c_k} \leq (A + Br^m)/r^k \to
0$ saat $r \to \infty$, sehingga $c_k = 0$: jadi $f =
\sum_{k\leq m}c_kz^k$ merupakan polinomial berderajat paling banyak $m$.

\textbf{13.} Fungsi $g = \eu^f$ utuh dengan $\abs g =
\eu^{\operatorname{Re}f} \leq \eu^M$: jadi konstan menurut pertanyaan
11. Lalu $0 = g' = f'g$ dengan $g$ yang bebas akar: sehingga $f' = 0$ dan $f$
konstan.

\textbf{14.} Misalkan $M = \sup_K\abs f$ pada bujur sangkar satuan
tertutup $K$ yang \omterm{def:b3:topology:compact}{kompak}. Setiap $z$ berselisih dari sebuah titik $K$ oleh sebuah
unsur $\Z + \iu\Z$ (kurangkan bagian bulatnya), dan
mengiterasikan kedua hubungan periodiknya membiarkan $f$
tak berubah: sehingga $\abs f \leq M$ pada seluruh $\C$, dan pertanyaan 11
membuat $f$ konstan. Jadi fungsi tak konstan yang invarian
terhadap kekisinya tak dapat utuh: sehingga fungsi eliptik
teori klasiknya harus membawa kutub --- yakni gerbang historis
menuju \cref{ch:b3:residues}.

\textbf{15.} Jika $f(\C)$ melewatkan cakram $D(a, r)$, maka
$\abs{f(z) - a} \geq r$ untuk setiap $z$, sehingga $g = 1/(f - a)$ bersifat
utuh dengan $\abs g \leq 1/r$: jadi konstan menurut Liouville, sehingga
$f$ konstan. Kontrapositifnya: jangkauan sebuah fungsi utuh
yang tak konstan memotong setiap cakram --- yakni padat di $\C$.

\textbf{16.} Pilihlah $R$ dengan $\abs f \geq 1$ di luar
$D(0, R)$. Maka akar $f$ terletak di $\bar D(0, R)$ yang
\omterm{def:b3:topology:compact}{kompak}; dan seandainya tak berhingga banyak, akarnya akan menumpuk
di sana, sehingga \cref{thm:b3:holomorphic:identity} akan memaksa $f
\equiv 0$ --- yang mustahil. Sebutlah akarnya $z_1, \dots, z_p$, dengan
kelipatan $m_1, \dots, m_p$, lalu tetapkan $M = \sum m_i$ dan
$\Pi(z) = \prod_i(z - z_i)^{m_i}$. Dengan memfaktorkan setiap akarnya keluar
dari deret pangkatnya, $g = f/\Pi$ bersifat utuh dan bebas akar.
Untuk $\abs z \geq \max(R, 2\max_i\abs{z_i})$: $\abs{z - z_i}
\leq 2\abs z$ dan $\abs f \geq 1$, sehingga $\abs{1/g} =
\abs\Pi/\abs f \leq 2^M\abs z^M$; sedangkan pada cakram \omterm{def:b3:topology:compact}{kompak}
sisanya $1/g$ \omterm{def:b3:topology:continuity}{kontinu}, jadi terbatas: sehingga $\abs{1/g} \leq A +
B\abs z^M$ di mana-mana. Menurut pertanyaan 12, $1/g$ merupakan
polinomial; dan ia bebas akar, sehingga menurut pertanyaan 7 ia
konstanta tak nol $c$: jadi $f = \frac1c\Pi$ merupakan polinomial.
Sebaliknya, pertanyaan 1 membuat setiap polinomial tak konstan
meledak di tak hingga. Sedangkan $\eu^z$ sungguh terkecualikan: sebab sepanjang $\R_-$,
$\abs{\eu^z} = \eu^x \to 0$ padahal $\abs z \to \infty$.

\textbf{17.} Parametrikanlah \cref{thm:b3:holomorphic:formula}
di pusatnya: dengan $w = a + r\eu^{\iu\theta}$ dan $\dd w = \iu
r\eu^{\iu\theta}\dd\theta$,
\[
f(a) = \frac1{2\iu\pi}\int_{C_r}\frac{f(w)}{w - a}\,\dd w
= \frac1{2\pi}\int_0^{2\pi}
f\bigl(a + r\eu^{\iu\theta}\bigr)\,\dd\theta ;
\]
lalu mengambil bagian realnya memberikan sifat nilai rata-rata $u$. Jika
$u$ mencapai maksimum di sebuah titik \omterm{def:b3:topology:interior}{interior} $\Omega$ yang
\omterm{def:b3:topology:connected}{terhubung}: maka $\abs{\eu^f} = \eu^u$ mencapai maksimum di \omterm{def:b3:topology:interior}{interiornya},
sehingga $\eu^f$ konstan menurut
\cref{thm:b3:holomorphic:maximum}(2), dan $u =
\log\abs{\eu^f}$ konstan. Sedangkan pada daerah terbatas dengan
\omterm{def:b3:topology:continuity}{kekontinuan} sampai batasnya, $\sup_{\bar\Omega}u =
\sup_{\partial\Omega}u$, tepat seperti bagi $\abs f$.

\textbf{18.} \emph{Kasus $\abs a > r$.} Pilihlah $R$ dengan $r < R
< \abs a$: maka pada cakram cembung $D(0, R)$ fungsi $a - z$
bersifat \omterm{def:b3:holomorphic:holo}{holomorfik} dan bebas akar, dan $z \mapsto -1/(a - z)$ mempunyai
antiturunan $L$ di sana (\cref{thm:b3:holomorphic:cauchy});
lalu setelah konstantanya disesuaikan, $\bigl(\eu^{-L}(a -
z)\bigr)' = \eu^{-L}\bigl(-L'\,(a - z) - 1\bigr) = 0$ memberikan
$\eu^L = a - z$: sehingga sebuah logaritma \omterm{def:b3:holomorphic:holo}{holomorfik} ada, dan
$\log\abs{a - z} = \operatorname{Re}L(z)$. Lalu sifat nilai rata-rata
pertanyaan 17 di $0$ dengan jari-jari $r$:
\[
\frac1{2\pi}\int_0^{2\pi}
\log\bigl|a - r\eu^{\iu\theta}\bigr|\,\dd\theta
= \operatorname{Re}L(0) = \log\abs a .
\]
\emph{Kasus $\abs a < r$.} Dari $a - r\eu^{\iu\theta} =
-r\eu^{\iu\theta}\bigl(1 - \frac ar\eu^{-\iu\theta}\bigr)$,
rata-ratanya sama dengan $\log r$ ditambah rata-rata $\log\abs{1 -
\frac ar\eu^{-\iu\theta}}$. Lalu substitusi $\theta \mapsto
2\pi - \theta$, kemudian --- dengan menulis $\frac ar =
\rho\eu^{\iu\varphi}$ dengan $\rho < 1$, sebab kasus $a = 0$
trivial --- pergeseran $\theta \mapsto \theta - \varphi$
(sebab keduanya mengawetkan rata-rata pada satu periode) mengubahnya menjadi rata-rata
$\log\abs{1 - \rho\eu^{\iu\theta}}$: yakni kasus pertamanya dengan
$(a, r) = (1, \rho)$, yang memberikan $\log 1 = 0$. Totalnya:
$\log r = \log\max(\abs a, r)$ pada kedua kasusnya.

\textbf{19.} Berlaku $\log\abs{P(\eu^{\iu\theta})} = \log\abs c +
\sum_i\log\abs{\alpha_i - \eu^{\iu\theta}}$, dengan setiap akarnya
berulang menurut kelipatannya; lalu merata-ratakannya terhadap $\theta$ dan
menerapkan pertanyaan 18 dengan $r = 1$ pada setiap akar di luar lingkaran
satuannya melahirkan $\log\bigl(\abs c\prod_i\max(1,
\abs{\alpha_i})\bigr)$. Untuk pemeriksaannya. Bagi $P = 2X - 1 = 2(X -
\frac12)$ rumusnya meramalkan $\log 2$; sedangkan secara langsung,
$\abs{2\eu^{\iu\theta} - 1} = 2\abs{\frac12 -
\eu^{\iu\theta}}$ dan rata-rata $\log\abs{\frac12 -
\eu^{\iu\theta}}$ adalah $\log\max(\frac12, 1) = 0$: jadi rata-ratanya $\log
2$. Bagi $P = X^2 - X = X(X - 1)$ akarnya $1$ terletak \emph{pada}
lingkarannya; dan rumusnya meramalkan $0$. Secara langsung, rata-rata
$\log\abs{\eu^{\iu\theta}}$ adalah $0$, dan dengan
$\abs{\eu^{\iu\theta} - 1} = 2\abs{\sin\frac\theta2}$:
\[
\frac1{2\pi}\int_0^{2\pi}
\log\Bigl(2\sin\frac\theta2\Bigr)\dd\theta
= \frac1\pi\int_0^\pi\log(2\sin u)\,\dd u
= \log 2 + \frac J\pi,
\qquad J = \int_0^\pi\log\sin u\,\dd u .
\]
Lalu dengan substitusi $u = 2v$, dan dengan memakai identitas
$\sin 2v = 2\sin v\cos v$, diperoleh
$J = \pi\log2 + 2\int_0^{\pi/2}\log\sin +
2\int_0^{\pi/2}\log\cos = \pi\log 2 + 2J$ (sebab masing-masing separuhnya sama dengan
$J/2$ menurut kesetangkupan $\sin$), sehingga $J = -\pi\log 2$ (sebab
integral tak wajarnya konvergen, karena $\log\sin$ \omterm{def:b3:lebesgue:l1}{terintegralkan} di
ujungnya): jadi rata-ratanya $\log2 - \log2 = 0$. Sehingga rumusnya
bertahan bagi akar pada lingkarannya.

\textbf{20.} Fungsi $\frac{\eu^z - 1}z =
\sum_{j\geq0}\frac{z^j}{(j+1)!}$ bersifat utuh dan bernilai $1$ di
$0$: sehingga kebalikannya \omterm{def:b3:holomorphic:holo}{holomorfik} di dekat $0$ (sungguh pada $\abs z <
2\pi$, sebab akar lain terdekat $\eu^z - 1$ adalah
$\pm2\iu\pi$), jadi $\frac z{\eu^z-1}$ analitik di $0$.
Lalu dengan mengalikan kedua deretnya dan membaca koefisien
$z^n$ dengan $n \geq 1$ pada $(\frac{\eu^z-1}z)\cdot(\frac
z{\eu^z-1}) = 1$:
\[
\sum_{k=0}^{n}\frac{B_k}{k!\,(n+1-k)!} = 0
\quad\Longleftrightarrow\quad
\sum_{k=0}^{n}\binom{n+1}{k}B_k = 0 .
\]
Berturut-turut: $B_0 = 1$, $B_1 = -\frac12$, $B_2 = \frac16$,
$B_3 = 0$, $B_4 = -\frac1{30}$, $B_5 = 0$, dan $B_6 =
\frac1{42}$. Untuk kegenapannya: dengan $F(z) = \frac z{\eu^z-1} +
\frac z2$,
\[
F(-z) = \frac{-z}{\eu^{-z} - 1} - \frac z2
= \frac{z\,\eu^z}{\eu^z - 1} - \frac z2
= z + \frac{z}{\eu^z - 1} - \frac z2 = F(z) :
\]
maka $F$ genap, sehingga $B_{2k+1} = 0$ untuk $k \geq 1$ (sebab satu-satunya
koefisien gasalnya $B_1$ terserap oleh $+\frac z2$). Untuk penunjuk
ke depannya: $\cot w = \iu + \frac{2\iu}{\eu^{2\iu w}-1}$ memberikan
$w\cot w = 1 + \sum_{k\geq1}\frac{B_{2k}}{(2k)!}(2\iu
w)^{2k}$, sehingga koefisien Laurent kotangennya ---
jadi, menurut \cref{ch:b3:residues}, setiap $\zeta(2k)$ --- dihargai
oleh bilangan Bernoulli:
\[
\zeta(2k) = (-1)^{k+1}\,\frac{(2\pi)^{2k}B_{2k}}{2\,(2k)!} .
\]

\textbf{21.} Tulis $f = \sum_nc_nz^n$ (dengan jari-jari $\infty$);
maka $g(z) = \overline{f(\bar z)} = \sum_n\bar c_nz^n$ bersifat
utuh. Pada $\R$: $g(x) = \overline{f(x)} = f(x)$, sehingga $g$ dan
$f$ bersesuaian pada himpunan yang bertitik tumpuk di $\C$ yang
\omterm{def:b3:topology:connected}{terhubung}: sehingga \cref{thm:b3:holomorphic:identity} memberikan $g \equiv f$,
yakni $f(\bar z) = \overline{f(z)}$ (setara dengan itu: semua
$c_n$ real). Untuk polinomial real $P$:
$P(\bar\alpha) = \overline{P(\alpha)} = 0$, dan identitas yang sama
yang diterapkan pada turunan realnya $P', P'', \dots$
mengawetkan kelipatannya: sehingga akar tak realnya berpasangan ---
yakni pemasangan pertanyaan 8, yang dibuktikan kembali secara analitik.

\textbf{22.} Tangganya, yang dirakit: terbatas $\Rightarrow$
konstan (11); terdominasi oleh $A + B\abs z^m$ $\Rightarrow$
polinomial (12); bagian real terbatas dari atas $\Rightarrow$
konstan (13); berperiode ganda $\Rightarrow$ konstan (14);
jangkauannya melewatkan sebuah cakram $\Rightarrow$ konstan (15); meledak di tak hingga
$\Rightarrow$ polinomial (16). Setiap anak tangganya adalah rumus
Cauchy: sebab nilai di sebuah titik merupakan rata-rata lingkaran, sehingga semua
koefisien Taylornya dihargai oleh \omterm{def:b3:measure:measure}{ukuran} $f$ pada lingkaran
yang besar, dan sebuah batas pertumbuhan memusnahkan koefisiennya
secara borongan. Tak ada yang seperti itu mengendalai fungsi $\mathcal
C^\infty$ real: sebab \omterm{def:b3:lp:mollifier}{fungsi gundukan}
(\cref{thm:b3:lp:regularization}) bersifat terbatas, berpendukung
\omterm{def:b3:topology:compact}{kompak}, dan liar tak konstan, sedangkan semua turunannya di sembarang
titik di luar pendukungnya lenyap tanpa fungsinya
lenyap di mana pun di dekatnya. Jadi kemulusan sama sekali tak menggandengkan turunannya
di titik yang berbeda; sedangkan \omterm{def:b3:holomorphic:holo}{keholomorfikan} merantai setiap
turunan ke satu integral atas lingkaran yang jauh. Yakni sebuah
syarat lokal dengan pemberi kabar global --- itulah mengapa
fungsi utuh menuruti hukum dan ketertiban.

\textbf{23.} Dengan menguraikan
$\abs{P(\eu^{\iu\theta})}^2 = \sum_{k,l}a_k\bar a_l
\eu^{\iu(k-l)\theta}$ lalu merata-ratakannya, setiap suku $k \neq
l$ terbunuh: sehingga rata-ratanya $\sum_k\abs{a_k}^2 =: N$. Andaikan dahulu bahwa
$P$ tak berakar pada lingkaran satuannya, sehingga $\theta \mapsto
\log\abs{P(\eu^{\iu\theta})}$ \omterm{def:b3:topology:continuity}{kontinu}. Lalu batas $\log
t \leq t - 1$ yang diterapkan pada $t = \abs P^2/N$ memberikan, setelah
dirata-ratakan,
\[
\frac1{2\pi}\int_0^{2\pi}\log\abs{P(\eu^{\iu\theta})}^2
\dd\theta - \log N
\;\leq\; \frac1N\cdot N - 1 = 0,
\]
sehingga rata-rata ukur $\abs P$ paling banyak $\sqrt N$;
lalu pertanyaan 19 mengenali rata-rata ukur itu sebagai
$\abs c\prod_i\max(1, \abs{\alpha_i})$: yakni ketaksamaan
Landau. Untuk akar pada lingkarannya: pilihlah $r > 1$ yang berbeda dari
setiap $\abs{\alpha_i}$; maka polinomial $P(rX)$, yang berakar
$\alpha_i/r$ di luar lingkaran satuannya dan berkoefisien $a_kr^k$,
memenuhi ketaksamaannya; lalu kedua ruasnya \omterm{def:b3:topology:continuity}{kontinu} terhadap $r$,
sehingga membiarkan $r \to 1^+$ memberikan kasus umumnya. Pada $X^2 -
X$: akarnya $0$ dan $1$, sehingga ruas kirinya $1$, sedangkan
ruas kanannya $\sqrt{1 + 1} = \sqrt2$: jadi benar, dengan kelonggaran.

\textbf{24.} Tulis $\eu^z - 1 = z\,g(z)$ dengan $g(z) =
\sum_{k\geq0}\frac{z^k}{(k+1)!}$ yang utuh dan $g(0) = 1$. Karena
$\eu^z = 1$ tepat pada $2\pi\iu\Z$, maka $g$ tak berakar di
$D(0, 2\pi)$ (untuk $0 < \abs z < 2\pi$ sebab $\eu^z - 1
\neq 0$, dan di $0$ sebab $g(0) = 1$), sehingga $h = 1/g$ \omterm{def:b3:holomorphic:holo}{holomorfik}
pada $D(0,2\pi)$ dan deret Taylornya di $0$ --- yang menurut
definisinya $\sum\frac{B_n}{n!}z^n$ --- konvergen pada
seluruh cakramnya: jadi $\rho \geq 2\pi$. Seandainya $\rho > 2\pi$, jumlahnya $S$
akan \omterm{def:b3:holomorphic:holo}{holomorfik} pada $D(0,\rho)$, dan ia bersesuaian dengan $z
\mapsto z/(\eu^z - 1)$ pada $0 < \abs z < 2\pi$; lalu keduanya
\omterm{def:b3:holomorphic:holo}{holomorfik} pada \omterm{def:b3:topology:topology}{himpunan terbuka} \omterm{def:b3:topology:connected}{terhubung} $D(0,\rho) \setminus
2\pi\iu\Z$, sehingga menurut teorema keidentikannya keduanya bersesuaian di sana. Padahal
saat $z \to 2\pi\iu$, $\abs{z/(\eu^z - 1)} \to \infty$
(sebab pembilangnya $\to 2\pi$ dan penyebutnya $\to 0$) sedangkan $S$
\omterm{def:b3:topology:continuity}{kontinu} di $2\pi\iu$: yang bertentangan. Karena itu $\rho = 2\pi$
secara persis, dan rumus Hadamard memberikan
$\limsup_n\abs{B_n/n!}^{1/n} = \frac1{2\pi}$; lalu karena koefisien
gasalnya nol mulai dari $n = 3$, $\limsup$-nya dibawa
\omterm{def:b3:holomorphic:index}{indeks} genapnya, yang merupakan rumus yang dinyatakan itu
dengan $n = 2k$. Secara numerik di $k = 6$: $(2\pi)^{12} \approx
3.7858\cdot10^9$ dan $2\cdot12! = 958\,003\,200$, sehingga
$2\,(2k)!/(2\pi)^{2k} \approx 0.25305$, terhadap
$\abs{B_{12}} = \frac{691}{2730} \approx 0.25311$. Sedangkan
nisbahnya, $1.00025$, tepat $\zeta(12)$ sampai angka yang
ditampilkan: sehingga rumus kalkulus residu $\zeta(2k) =
(-1)^{k+1}\frac{(2\pi)^{2k}B_{2k}}{2\,(2k)!}$ pada
\cref{ch:b3:residues} menerangkan baik faktor $2$-nya maupun
kelebihan mungilnya.

\textbf{25.} Untuk batas Cauchynya: jika $\abs z > 1 + M$ dengan $M =
\max_k\abs{q_k}$, maka
\[
\Bigl|\sum_{k<n}q_kz^k\Bigr|
\leq M\,\frac{\abs z^n - 1}{\abs z - 1}
< \frac{M}{\abs z - 1}\,\abs z^n \leq \abs z^n,
\]
sehingga $\abs{Q(z)} > 0$: jadi semua akarnya terletak di $\overline D(0, 1+M)$.
Lalu koefisien $P_j$ konvergen, sehingga terbatas oleh
suatu $M$: jadi semua akar semua $P_j$ (dan akar $P$) terletak di
$K = \overline D(0, 1 + M)$ yang \omterm{def:b3:topology:compact}{kompak}. Misalkan $v_j \in K^n$ sebuah
vektor yang mendaftarkan akar $P_j$ beserta kelipatannya
(pertanyaan 7). Maka setiap subbarisan $(v_j)$ mempunyai subbarisan
lanjutan yang konvergen ke suatu $(\beta_1, \dots, \beta_n)$;
sedangkan koefisien $\prod_i(X - \alpha_i^{(j)})$ adalah, sampai
tandanya, fungsi setangkup elementer $v_j$ ---
yang \omterm{def:b3:topology:continuity}{kontinu} --- sehingga sepanjang subbarisan itu keduanya konvergen ke
koefisien $\prod_i(X - \beta_i)$; padahal keduanya
konvergen ke koefisien $P$ menurut hipotesisnya, sehingga $\prod_i(X -
\beta_i) = P$: jadi setiap limit subbarisan $(v_j)$ merupakan
permutasi vektor akar $P$. Seandainya jarak pemasangannya
$\delta_j = \min_\sigma\max_i\,
\abs{\alpha_i^{(j)} - \alpha_{\sigma(i)}}$ tak menuju
$0$, maka sebuah subbarisan akan menjaga $\delta_j \geq \varepsilon$
sedangkan vektor akarnya konvergen ke sebuah permutasi
akar $P$ --- yang memaksa $\delta_j \to 0$ sepanjangnya:
yang bertentangan. Jadi multihimpunan akarnya konvergen. Untuk ketajamannya:
$P_\varepsilon = (X - 1)^2 + \varepsilon$ berakar $1 \pm
\iu\sqrt\varepsilon$: sehingga akar rangkapnya bergeser sejauh
$\sqrt\varepsilon$, misalnya sejauh $10^{-2}$ untuk $\varepsilon =
10^{-4}$. Secara umum, jika $\alpha$ akar berlipat $m$, maka
di dekat $\alpha$ berlaku $\abs{P(z)} \asymp \abs{z -
\alpha}^m$, sehingga gangguan berukuran $\varepsilon$ menggeser
gugusan akarnya kira-kira sejauh $\varepsilon^{1/m}$: yakni \omterm{def:b3:topology:continuity}{kekontinuan}
Hölder bereksponen $\frac1m$ dan tak lebih baik --- itulah
mengapa penyelesai numerik di dekat akar rangkap hanya menyimpan separuh
angka kerjanya.

\end{solution}
